\documentclass[11pt,a4paper,fontset=fandol]{ctexart}

\usepackage[a4paper,top=2.25cm,bottom=2.45cm,left=2.25cm,right=2.25cm,headheight=18pt,headsep=0.55cm,footskip=24pt]{geometry}
\usepackage{amsmath,amssymb,amsthm}
\usepackage{fancyhdr}
\usepackage{lastpage}
\usepackage{enumitem}
\usepackage{titlesec}
\usepackage{xcolor}
\usepackage{tikz}
\usepackage{hyperref}
\usepackage{bookmark}
\usepackage[most]{tcolorbox}
\usepackage{setspace}
\usepackage{array}
\usetikzlibrary{positioning}

\hypersetup{
  colorlinks=true,
  linkcolor=blue!50!black,
  urlcolor=blue!50!black
}

\setstretch{1.13}
\setlength{\parindent}{2em}
\setlength{\parskip}{0.25em}
\allowdisplaybreaks
\setlist[enumerate]{itemsep=0.45em, topsep=0.35em, leftmargin=2.4em}
\setlist[itemize]{itemsep=0.25em, topsep=0.25em, leftmargin=1.9em}

\definecolor{mainblue}{RGB}{36,83,140}
\definecolor{lightblue}{RGB}{241,246,252}
\definecolor{lightgray}{RGB}{246,246,246}
\definecolor{accent}{RGB}{153,76,0}
\definecolor{rulegray}{RGB}{110,118,128}

\titleformat{\section}
  {\Large\bfseries\color{mainblue}}
  {\thesection.}
  {0.45em}
  {}
  [\vspace{0.2em}\color{rulegray}\titlerule]
\titlespacing*{\section}{0pt}{1.1em}{0.7em}

\pagestyle{fancy}
\fancyhf{}
\fancyhead[L]{\small 图论计数周测 A 卷}
\fancyhead[C]{\small 组合专题：握手定理与基础建图}
\fancyhead[R]{\small 刘文博}
\fancyfoot[C]{\small 第 \thepage 页 / 共 \pageref{LastPage} 页}
\renewcommand{\headrulewidth}{0.55pt}
\renewcommand{\footrulewidth}{0.35pt}

\newtcolorbox{infobox}[1]{
  breakable,
  colback=lightblue,
  colframe=mainblue,
  boxrule=0.7pt,
  arc=1mm,
  title=\textbf{#1},
  fonttitle=\bfseries
}

\newenvironment{solution}{\par\noindent\textbf{参考解答：}}{\hfill$\square$\par}
\newcommand{\answerspace}{\vspace{2.3cm}}
\newcommand{\biganswerspace}{\vspace{3.1cm}}

\begin{document}

\thispagestyle{empty}
\begin{center}
{\fontsize{22}{28}\selectfont\bfseries 图论计数周测 A 卷\par}
\vspace{0.4cm}
{\large 适合小学高年级至初中低年级数学竞赛训练\par}
\end{center}

\begin{infobox}{考试说明}
\begin{itemize}
  \item 测试时间：70--75 分钟
  \item 满分：100 分
  \item A 卷侧重点、边、度数、完全图、二分图和基础网格路径，建议先把题目翻译成图再下手。
\end{itemize}
\end{infobox}

\section{试题}

\begin{enumerate}
  \item （10 分）完全图 $K_8$ 有多少条边？
  \answerspace

  \item （12 分）一个图有 10 个点，每个点的度数都等于 4。求这个图的边数。
  \answerspace

  \item （12 分）证明：图中不可能恰好有 3 个奇度数点。
  \answerspace

  \item （12 分）一个二分图两部分点数分别为 6 和 7，问边数最多是多少？
  \begin{center}
  \begin{tikzpicture}[
    scale=0.85,
    dot/.style={circle,draw=mainblue,fill=lightblue,inner sep=1.3pt},
    edge/.style={draw=black!60,line width=0.8pt}
  ]
    \foreach \y in {1.5,0.5,-0.5,-1.5}
      \node[dot] at (0,\y) {};
    \foreach \y in {1.8,1.2,0.6,0,-0.6,-1.2,-1.8}
      \node[dot] at (4,\y) {};
    \draw[edge] (0,1.5)--(4,1.8);
    \draw[edge] (0,0.5)--(4,0.6);
    \draw[edge] (0,-0.5)--(4,-0.6);
    \draw[edge] (0,-1.5)--(4,-1.8);
    \node[draw=none,fill=none,circle=false,font=\small] at (0,-2.4) {一侧 6 点};
    \node[draw=none,fill=none,circle=false,font=\small] at (4,-2.4) {另一侧 7 点};
  \end{tikzpicture}
  \end{center}
  \answerspace

  \item （12 分）8 个同学参加循环赛，每两人比赛 1 场，共有多少场比赛？
  \answerspace

  \item （14 分）从 $(0,0)$ 到 $(4,3)$，每次只能向右或向上走 1 格，共有多少条最短路径？
  \answerspace

  \item （14 分）从 $(0,0)$ 到 $(4,4)$，每次只能向右或向上走 1 格，但要求必须经过点 $(2,1)$。共有多少条最短路径？
  \begin{center}
  \begin{tikzpicture}[scale=0.62]
    \draw[step=1cm,gray!55,thin] (0,0) grid (4,4);
    \fill[mainblue] (0,0) circle (2pt);
    \fill[mainblue] (4,4) circle (2pt);
    \fill[accent!85!black] (2,1) circle (2.5pt);
    \draw[accent!85!black,line width=0.9pt] (2,1) circle (0.18);
    \node[below left] at (0,0) {\small $(0,0)$};
    \node[above right] at (4,4) {\small $(4,4)$};
    \node[below] at (2,1) {\small 必经点 $(2,1)$};
  \end{tikzpicture}
  \end{center}
  \biganswerspace

  \item （14 分）某个聚会共有 9 人，每个人都恰好认识 4 人。问一共有多少对认识关系？
  \biganswerspace
\end{enumerate}

\newpage
\section{详细解答}

\begin{enumerate}
  \item \begin{solution}
  完全图 $K_8$ 的边数就是从 8 个点中任选 2 个点的方法数：
  \[
  \binom82=28.
  \]
  \end{solution}

  \item \begin{solution}
  总度数为
  \[
  10\times 4=40.
  \]

  由握手定理
  \[
  2E=40,
  \]
  所以
  \[
  E=20.
  \]
  \end{solution}

  \item \begin{solution}
  由握手定理，图中所有点的度数之和一定是偶数。

  如果恰好有 3 个奇度数点，那么奇度数点的度数和是奇数；其余偶度数点的度数和是偶数；奇数加偶数仍是奇数。

  这就与总度数和必须是偶数矛盾。

  因此图中不可能恰好有 3 个奇度数点。
  \end{solution}

  \item \begin{solution}
  二分图两部分点数分别为 6 和 7 时，边数最多为
  \[
  6\times 7=42.
  \]
  \end{solution}

  \item \begin{solution}
  8 个同学两两比赛 1 场，就是完全图 $K_8$ 的边数，因此总场数为
  \[
  \binom82=28.
  \]
  \end{solution}

  \item \begin{solution}
  从 $(0,0)$ 到 $(4,3)$，要走 4 步向右和 3 步向上，共 7 步。

  所以最短路径条数为
  \[
  \binom73=35.
  \]
  \end{solution}

  \item \begin{solution}
  先从 $(0,0)$ 走到 $(2,1)$：需要 2 步向右、1 步向上，所以有
  \[
  \binom31=3
  \]
  条。

  再从 $(2,1)$ 走到 $(4,4)$：需要 2 步向右、3 步向上，所以有
  \[
  \binom53=10
  \]
  条。

  因此经过 $(2,1)$ 的最短路径共有
  \[
  3\times 10=30
  \]
  条。
  \end{solution}

  \item \begin{solution}
  把每个人看成一个点，把“认识”看成一条无向边。

  因为每个人都恰好认识 4 人，所以总度数是
  \[
  9\times 4=36.
  \]

  由握手定理
  \[
  2E=36,
  \]
  所以
  \[
  E=18.
  \]

  因此一共有
  \[
  18
  \]
  对认识关系。
  \end{solution}
\end{enumerate}

\end{document}
